% **************************************************************************************************************
%
%  Copyright 2020-2026 Robert Bosch GmbH
%
%  Licensed under the Apache License, Version 2.0 (the "License");
%  you may not use this file except in compliance with the License.
%  You may obtain a copy of the License at
%
%      http://www.apache.org/licenses/LICENSE-2.0
%
%  Unless required by applicable law or agreed to in writing, software
%  distributed under the License is distributed on an "AS IS" BASIS,
%  WITHOUT WARRANTIES OR CONDITIONS OF ANY KIND, either express or implied.
%  See the License for the specific language governing permissions and
%  limitations under the License.
%
% **************************************************************************************************************

\section{History}

\subsection{Version 2.1.0 - 11.05.2026}

First minor release after the gRPC + Consul + Nomad migration shipped
in 2.0.0.  Covers the Manager-GUI parity work, the Service Creator
wizard upgrades, the installer modernisation, and the new Robot
Framework test path.

\subsubsection{Architecture migration to gRPC + Consul + Nomad}

\begin{itemize}
   \item \textbf{Transport}: RabbitMQ (custom request/response framing) replaced by
         gRPC over HTTP/2 with protobuf payloads.
   \item \textbf{Discovery}: in-process \texttt{ServiceRegistry} replaced by the
         Consul service catalog (\texttt{/v1/health/service/...}).
   \item \textbf{Orchestration}: custom Local Process Hub replaced by Nomad jobs
         (\texttt{raw\_exec} driver) supervised by a Nomad agent.
   \item \textbf{Routing}: routing keys + alias routing replaced by
         fully-qualified gRPC service + method names.
   \item \textbf{Topology}: per-service \texttt{deploy/<svc>.nomad.hcl} replaces the
         single \texttt{hub\_processes.json} config.
   \item Hexagonal architecture itself is unchanged --- the migration swapped
         adapters, not domain code.
\end{itemize}

\subsubsection{New Architecture Decision Records (ADRs 023--029)}

\begin{itemize}
   \item \textbf{ADR-023}: Multi-node Consul cluster with serf gossip.
   \item \textbf{ADR-024}: Multi-node Nomad cluster (server + N clients).
   \item \textbf{ADR-025}: Wrapper layer for Nomad \texttt{raw\_exec} workloads
         (PATH setup, env-var derivation, pre-flight checks).
   \item \textbf{ADR-026}: \texttt{uv}-based Python service environments
         (per-service venv, lockfile).
   \item \textbf{ADR-027}: Kafka event bus alongside gRPC (sync = gRPC,
         fan-out = Kafka).
   \item \textbf{ADR-028}: gRPC + reflection as the canonical RPC layer
         (consolidates the rationale scattered across the ADR-016/017/018
         supersession callouts).
   \item \textbf{ADR-029}: Robot Framework as the primary test client, via the
         new \texttt{GrpcClient} connection type for QConnectBase.
   \item Older ADRs (003, 010--012, 015--018, 020) carry visible
         \emph{Superseded} or \emph{Archived} callout banners pointing at the
         new ones.
\end{itemize}

\subsubsection{Service Creator wizard}

\begin{itemize}
   \item \textbf{Multi-proto layout}: one process hosting N gRPC services on
         the same \texttt{ServerBuilder} / \texttt{grpc.aio.Server}, one Consul
         registration, one port.  Best for one logical device with multiple
         capability surfaces.
   \item \textbf{Multi-file \texttt{.proto} import}: select N \texttt{.proto}
         files at once; every service from every file gets loaded; layout
         auto-defaults to \texttt{multi\_proto}.
   \item \textbf{Three layout options} on the picker modal: \texttt{multi\_proto},
         \texttt{monorepo} (one project, N entry points), \texttt{separate}
         (N independent project folders).  Inline layout selector also added on
         Step 3 so users can flip without re-importing.
   \item \textbf{New ``GUI Support'' wizard step} (Python + HTML/JS only) with
         two tabs: a Schema Builder that emits \texttt{gui\_schema.json} (4
         component types: Method Form, Result Table, Static Text, Live Status)
         and a Custom HTML/JS editor with live preview and JS file upload.
   \item \textbf{Server toolchain selector} on Step 2 (MSYS2 prebuilt vs.
         vcpkg + Qt MinGW) independent of the client gRPC stack choice; mixed
         configurations get inline warnings.
   \item \textbf{Pre-generate proto stubs} on the bridge runs in-process via
         \texttt{grpc\_tools.protoc}; failures now log clearly and the GUI
         surfaces a yellow toast when stubs were requested but missing.
   \item \textbf{Language-aware copy}: \texttt{.exe} references replaced by
         language-appropriate wording (``Python entry point'' for Python,
         \texttt{.exe} for C++).
\end{itemize}

\subsubsection{Generated services: full documentation}

\begin{itemize}
   \item Generated Python and C++ services now ship with QConnect-style
         file headers (Apache copyright, File / Author / Date / Description
         / History sections).  Author is the OS user's full display name
         (\texttt{GetUserNameExW} on Windows, \texttt{pwd.pw\_gecos} on Linux);
         date and history are stamped at generation time.
   \item Per-method docstrings (Python: QConnect format with
         \textbf{Arguments:} / \textbf{Returns:} blocks; C++: Doxygen
         \texttt{@param} / \texttt{@return} blocks).
   \item Class docstrings on Settings, Context, Domain, GrpcAdapter.
   \item TODO-stub bodies in domain methods describe what the user should
         document and that the default returns the literal string
         \texttt{"not implemented"}.
\end{itemize}

\subsubsection{Manager GUI}

\begin{itemize}
   \item \textbf{Settings dialog}: new \emph{Service infrastructure} section
         with Consul / Nomad executable path fields and OS file-picker
         Browse buttons (Electron only).  Lookup priority documented inline:
         this setting $\rightarrow$ \texttt{PATH} $\rightarrow$ default install location.
   \item \textbf{Service Network mode}: Consul + Nomad dashboards (replaces
         the legacy Fleet view).  Consul tab can start / stop / view-log a
         managed dev-mode agent; same for Nomad.
   \item \textbf{Bridge endpoints} added: \texttt{/api/consul/agent/start} +
         \texttt{stop} + \texttt{status} + \texttt{log}; same for
         \texttt{/api/nomad/agent/*}.  Each persists the spawned agent's PID
         to a writable data dir for cross-restart reconnection.
   \item \textbf{InfraStatus pills} in the navbar (Consul / Nomad / Bridge)
         poll every 5 seconds; click jumps to the matching Service Network
         tab.
   \item \textbf{README + sidebar nav refresh}: documentation map moved from
         the top of the README to the bottom (less intrusive for first-time
         readers); per-page sidebar TOCs added.
\end{itemize}

\subsubsection{Installer modernisation (Windows + Linux)}

\begin{itemize}
   \item \textbf{NSIS}: new \emph{Infrastructure} wizard page detects Consul /
         Nomad on \texttt{PATH}, in \texttt{\%ProgramFiles64\%\textbackslash HashiCorp\textbackslash},
         in \texttt{\%LOCALAPPDATA\%\textbackslash Microsoft\textbackslash WinGet\textbackslash Links\textbackslash},
         and in \texttt{WinGet\textbackslash Packages\textbackslash Hashicorp.*\textbackslash}
         (\texttt{FindFirst} glob with named labels --- \texttt{+N} relative jumps
         and \texttt{\$PROGRAMDATA} / \texttt{\$LOCALAPPDATA} replaced with
         \texttt{ReadEnvStr} for portability across NSIS versions).
   \item \textbf{Three install options} per tool: install via winget (pinned
         versions \texttt{Consul 1.20.6} / \texttt{Nomad 1.9.6}), provide an
         existing path, or skip.  Winget calls use \texttt{--source winget} +
         \texttt{--silent --accept-*-agreements --disable-interactivity};
         re-detect-on-failure handles ``already installed'' as success.
   \item \textbf{Settings-JSON integration}: chosen Consul / Nomad paths
         written to \texttt{settings.json} so the GUI's Service Network tab
         can launch the agents without further configuration.
   \item \textbf{Legacy install steps deactivated}: Erlang/OTP, RabbitMQ
         Server, and ProcessHub install flows wrapped in
         \texttt{ENABLE\_LEGACY\_BROKER} / \texttt{ENABLE\_LEGACY\_PROCESSHUB}
         conditional-compile guards (default off; flip to 1 to re-enable).
   \item \textbf{Linux}: new \texttt{bootstrap.sh} interactive script shipped
         in \texttt{extraResources} for AppImage users, plus a \texttt{.deb}
         postinst hook that launches it when stdin is a TTY.  Detects distro,
         installs from HashiCorp's official apt / yum repos with version
         pinning, falls back to release-tarball download.
   \item \textbf{Installer size}: cut from \texttildelow 259 MB to \texttildelow 30--40 MB
         compressed by excluding \texttt{examples/\_legacy/} (\texttildelow 393 MB
         unpacked of pre-migration RabbitMQ-era templates), trimming
         \texttt{dist-msys2/} and \texttt{deploy/} build artefacts, and
         restricting Chromium locales to \texttt{en-US}.
\end{itemize}

\subsubsection{Robot Framework test path: QConnectBase \texttt{GrpcClient}}

\begin{itemize}
   \item New \texttt{QConnectBase/grpc/grpc\_client.py} contributed upstream
         to QConnectBase.  Subclasses \texttt{ConnectionBase} and integrates
         with the standard \texttt{Connect} / \texttt{Send Command} /
         \texttt{Verify} / \texttt{Disconnect} keyword library.
   \item Two addressing modes: direct (\texttt{target: host:port}) and
         Consul-resolved (\texttt{service\_name + consul\_addr}, lazy import
         of \texttt{MicroserviceBase.runtime.consul.resolve\_via\_consul}).
   \item \texttt{wait\_4\_trace} overridden to ignore the regex
         \texttt{search\_obj} (gRPC returns structured objects, not text
         streams) and to return the response as a Python \texttt{dict}.
         Tests assert content with Robot's built-in
         \texttt{Should Be Equal} / \texttt{Dictionary Should Contain}.
   \item Reflection client (\texttt{GrpcReflectClient}) with automatic
         fallback to \texttt{LocalProtoClient} when the server lacks
         \texttt{grpc++\_reflection} (e.g. C++ services built against the
         vcpkg grpc port).
\end{itemize}

\subsubsection{C++ runtime as standalone CMake / vcpkg package}

\begin{itemize}
   \item \texttt{runtime\_cpp/CMakeLists.txt} now exports
         \texttt{MicroserviceBaseConfig.cmake} +
         \texttt{MicroserviceBaseTargets.cmake} via
         \texttt{install(TARGETS ... EXPORT)}.  After
         \texttt{cmake -{}-install}, generated services can
         \texttt{find\_package(MicroserviceBase CONFIG REQUIRED)} from anywhere.
   \item New vcpkg overlay-port at \texttt{ports/microservice-base/}.
   \item Generator emits a hybrid CMake block: tries
         \texttt{find\_package} first, falls back to
         \texttt{add\_subdirectory(.../runtime\_cpp)} when the path resolves
         (i.e.\ the service lives inside \texttt{<framework>/examples/}).
   \item Three target names exported for back-compat:
         \texttt{microservice\_base::runtime} (modern alias),
         \texttt{microservice\_base::microservice\_base\_runtime} (vcpkg style),
         \texttt{microservice\_base\_runtime} (legacy bare name).
\end{itemize}

\subsubsection{Python monorepo and multi-proto generators}

\begin{itemize}
   \item \texttt{python\_tmpl.generate\_monorepo()} mirrors the C++ monorepo
         emitter: one project, N service modules sharing a single \texttt{.proto},
         one \texttt{pyproject.toml} registering each service as a
         \texttt{[project.scripts]} entry point.
   \item Service Creator wizard's monorepo radio is now valid for both
         Python and C++.
\end{itemize}

\subsubsection{Documentation refresh}

\begin{itemize}
   \item New canonical
         \texttt{docs/diagrams/00\_canonical\_architecture.puml} (cluster
         topology) and \texttt{component.puml} (per-workload component layout),
         both adapted from the TA reference architecture.
   \item Class diagrams (\texttt{class\_domain}, \texttt{class\_ports},
         \texttt{class\_adapters}) rewritten to reflect the post-migration
         framework boundary.
   \item Sequence diagrams \texttt{sequence\_communication},
         \texttt{sequence\_rpc}, \texttt{sequence\_registration},
         \texttt{sequence\_shutdown} redrawn for the gRPC + Consul + Nomad path.
   \item 14 pre-migration diagrams archived under
         \texttt{docs/diagrams/\_archive/} with a mapping README pointing to
         their successors.
   \item Service Creator wizard documentation rewritten end-to-end (Markdown
         + HTML mirror) to match the current 5-step wizard, with new
         screenshots for Step 2 (server toolchain), the picker modal (3 layout
         options), Step 5 review (layout badge), and the new GUI Support tab.
   \item Full audit pass on framework + scaffold + Tier-A source files
         brought docstrings into the QConnect house style (column-0 content,
         \textbf{Arguments:} / \textbf{Returns:} blocks).
   \item All 24 live PUML diagrams re-rendered with PlantUML 1.2021.01;
         16 stale PNGs archived under \texttt{docs/imgs/\_archive/}.
\end{itemize}

\subsubsection{Dependency change}

\begin{itemize}
   \item \texttt{grpcio-tools} promoted from the \texttt{[proto-tools]}
         optional extra into the base \texttt{dependencies} list.  The
         Service Creator wizard's \emph{Pre-generate proto stubs} feature
         needs it at runtime --- it was previously a silent no-op when the
         bridge's Python lacked it.  Cost: \texttildelow 30 MB on every install.
   \item \texttt{[project.optional-dependencies]} \texttt{proto-tools}
         block removed; \texttt{all} extra updated accordingly.
\end{itemize}

\subsection{Version 2.0.0 - 09.02.2026}

Major restructure: hexagonal architecture, dual-host GUI, multi-broker support,
local process hub, standalone installer.

\subsubsection{Architecture}

\begin{itemize}
   \item \textbf{Hexagonal Architecture}: Reorganized codebase into domain, ports, and adapters layers
   \item \textbf{Factory Pattern}: All components created through \texttt{factory.py} for easy swapping
   \item \textbf{Port Interfaces}: \texttt{TransportPort}, \texttt{RegistryPort}, \texttt{UIBridgePort}
   \item \textbf{Adapter Implementations}: RabbitMQ transport, RabbitMQ registry, FastAPI bridge,
         Local Hub Manager, Service Executor
\end{itemize}

\subsubsection{GUI v2.0}

\begin{itemize}
   \item \textbf{Dual Hosting}: Electron desktop app and browser mode via FastAPI bridge
   \item \textbf{Pure Web Frontend}: HTML/CSS/JS with Bootstrap 5 CDN (no Node.js dependencies)
   \item \textbf{Namespace}: \texttt{window.MicroserviceManager} (alias \texttt{MM}) replaces
         Node.js \texttt{global} and \texttt{require()} patterns
   \item \textbf{Service Plugins}: Dynamically loaded per-service UI components in \texttt{web/services/}
   \item \textbf{Dark Sidebar Theme}: Custom Bootstrap dark sidebar with accordion navigation
   \item \textbf{Login Modal}: Bootstrap modal replaces separate Electron popup window
   \item \textbf{Toast Notifications}: \texttt{MM.showToast()} replaces \texttt{notifier.notify()}
         and \texttt{dialog.showMessageBox()}
\end{itemize}

\subsubsection{Multi-Broker Support}

\begin{itemize}
   \item \textbf{Connection Map}: \texttt{MM.connections} tracks per-broker state
   \item \textbf{Reverse Lookup}: \texttt{serviceToBroker} maps service names to broker endpoints
   \item \textbf{Progressive Disclosure}: Broker headers shown only with multiple connections
   \item \textbf{Session Persistence}: Connection state saved via \texttt{sessionStorage}
\end{itemize}

\subsubsection{Local Process Hub}

\begin{itemize}
   \item \textbf{Hub Manager}: Start, stop, and monitor local service processes
   \item \textbf{Service Executor}: Process execution with \texttt{CTRL\_BREAK\_EVENT} for
         graceful shutdown on Windows
   \item \textbf{Config Persistence}: \texttt{hub\_processes.json} with \texttt{\$\{python\}} and
         \texttt{\$\{config\_dir\}} placeholders preserved across save operations
   \item \textbf{Hub Dashboard}: Real-time process status display in the GUI
   \item \textbf{REST API}: Full hub management via \texttt{/api/local-hub/*} endpoints
\end{itemize}

\subsubsection{Service Management}

\begin{itemize}
   \item \textbf{Service Import}: Import microservice packages from folders or zip archives
   \item \textbf{Service Creator}: Interactive wizard for scaffolding new microservices
   \item \textbf{Registry Shutdown}: \texttt{\_\_registry\_shutdown\_\_} sentinel notifies subscribers
         when a service unregisters
\end{itemize}

\subsubsection{Electron Packaging}

\begin{itemize}
   \item \textbf{Standalone Installer}: NSIS installer for Windows (\texttt{build.bat}),
         Linux packages (\texttt{build.sh})
   \item \textbf{Read-only Resources}: Scripts in \texttt{resources/python/}, app in \texttt{app.asar}
   \item \textbf{Writable User Data}: Settings, config, services, logs in
         \texttt{\%APPDATA\%/devatservgui/}
   \item \textbf{First-Run Init}: Template files copied from resources on first launch
   \item \textbf{Bridge Lifecycle}: Detached bridge process managed via PID file
\end{itemize}

\subsubsection{Windows Process Fixes}

\begin{itemize}
   \item \texttt{SIGBREAK} handler converted to \texttt{KeyboardInterrupt} for graceful Python cleanup
   \item Pika \texttt{start\_consuming()} replaced with \texttt{process\_data\_events(time\_limit=1)}
         loop for interruptible consumption
   \item \texttt{subprocess.PIPE} blocking prevented by using \texttt{FileHandler} instead of
         \texttt{StreamHandler} for logging in subprocess mode
\end{itemize}

\subsubsection{Documentation}

\begin{itemize}
   \item 12 Architecture Decision Records (ADR-001 through ADR-012)
   \item 12 PlantUML diagrams (overview, architecture, component, class, sequence, state)
   \item Updated README with architecture diagrams, quick start guide, and API reference
   \item Package documentation with Introduction, Description, and History
\end{itemize}

\subsection{Version 0.1.0 - 02/2023}

Initial version.

\begin{itemize}
   \item Basic microservice framework with RabbitMQ communication
   \item Service registration and discovery
   \item Electron-based GUI for service monitoring
\end{itemize}
